\section{A power saving for transverse inputs}
\label{sec:transverse}

The first step is a rigidity statement for almost extremizers of the
fractal estimate.  A hypothetical extremizing configuration produces a
probability distribution on points and incident tubes.  Information at
successive spatial scales first forces three identical one-dimensional
profiles, then forces those profiles to alternate only finitely often
between slopes zero and one.  Polynomial partitioning rules out the
remaining transition.

\begin{proposition}[Transverse power saving]\label{prop:transverse}
Fix a bounded frequency region.  There are $c>0$ and $L_0<\infty$,
depending only on that region, for which the following configuration
cannot exist when $L\ge L_0$.  Let $\mathcal Q$ be unit lattice cubes
inside a ball of radius $L^2$ in $\R^3$, with
\[
 \#\mathcal Q\ge L^{4-c},\qquad
 \#\{Q\in\mathcal Q:Q\cap B(z,r)\ne\varnothing\}
       \le L^c r^2\quad(r\ge1),
\]
where cube counts may equivalently use their centers, up to fixed
constants.  Suppose $v_1,v_2,v_3$ have frequency support in the
fixed region, $\|v_i\|_2\le1$, and
\[
 \sup_Q|Ev_i|\ge L^{-4/3-c}
 \qquad(Q\in\mathcal Q,\ i=1,2,3).
\]
Then it is impossible that every triple of unit frequency directions,
one from each support, has absolute determinant at least $L^{-c}$.
\end{proposition}

We argue by contradiction along $L\to\infty$ and $c\downarrow0$.
Every $o(1)$ in this section refers to this sequence or a subsequence.
A diagonal choice converts estimates valid with every fixed positive
power loss into $L^{o(1)}$ estimates; constants are fixed before making
that choice.  All rounding, support enlargements, and packet padding
are included in these subpower errors.  The same convention applies
when a scale is a fixed positive power of $L$.

\subsection{Uniform multilinear estimates and the tube model}

We need the multilinear estimates when transversality itself deteriorates
by a subpower.  An unspecified transversality-dependent constant in the
restriction theorem would not justify this use.  The following finite
induction supplies the required uniformity.

\begin{lemma}[Multilinear estimates with subpower transversality]
\label{lem:transverse-product}
Suppose three exact extension inputs have bounded frequency supports,
and every triple of their unit directions has determinant at least
$\nu$, where $0<\nu\le1$.  For each $\epsilon>0$ there are $C_\epsilon,M_\epsilon<\infty$
such that, on every ball of radius $S\ge2$,
\begin{equation}\label{eq:transverse-product}
 \int_{B_S}\prod_{i=1}^3|Ev_i(z)|\,dz
 \le C_\epsilon\nu^{-M_\epsilon}S^\epsilon
                         \prod_{i=1}^3\|v_i\|_2.
\end{equation}
The same conclusion holds for the sum over unit lattice cubes in the
ball of $\prod_i\sup_Q|Ev_i|$.  In particular, when $\nu=S^{-o(1)}$
both bounds are $S^{o(1)}\prod_i\|v_i\|_2$.
For weighted tubes of width $\Delta$ in a bounded region, the
corresponding multilinear Kakeya bound is
\begin{equation}\label{eq:transverse-kakeya}
 \int\prod_{i=1}^3
       \left(\sum_T w_i(T)\mathbf1_T\right)^{1/2}
 \le C_\epsilon\nu^{-M_\epsilon}\Delta^{3-\epsilon}
                  \prod_{i=1}^3\left(\sum_Tw_i(T)\right)^{1/2}.
\end{equation}
Here fixed enlargements of tubes and of the ambient region are allowed;
if those enlargements are subpower on a sequence, they cost a subpower.
\end{lemma}

\begin{proof}
Partition each direction support into neighborhoods of radius a sufficiently
small fixed power of $\nu$.  There are at most $\nu^{-O(1)}$ such
neighborhoods.  For each triple choose representative directions and send
them to the coordinate axes by a linear map.  Its norm, inverse norm,
and volume distortion are bounded by fixed powers of $\nu^{-1}$.
The transformed direction neighborhoods lie in fixed small neighborhoods
of the axes.  The weighted multilinear Kakeya inequality of
\cite[Theorem~1.15 and Remark~1.10]{BCT06}, applied there and summed over triples, proves
\eqref{eq:transverse-kakeya}.  Weights follow first by repetitions and
then by approximation.  Widths larger than a fixed constant can instead
be treated by the trivial bound, after increasing the inverse-power
constant.

For completeness, use this Kakeya inequality to obtain restriction
uniformity.  Let $A_B(S,\nu)$ be the least normalized product-integral
constant for inputs supported in $[-B,B]^2$.  Fix padding exponent
$\theta>0$ and Kakeya loss $\sigma>0$.  Cover a ball of radius $S$
by cubes of side comparable to $\sqrt S$.  The exact packet
decomposition restricts each input on a smaller cube to packets
incident to its $S^\theta\sqrt S$ enlargement.  Its retained exact
input has squared norm at most a fixed constant times the incident
coefficient energy.  Its support is within $O_B(S^{-1/2})$ of the
original support: only caps intersecting that support occur, and
every enlarged cap lies within this distance of such an intersection.
Thus determinant continuity preserves $\nu/2$ transversality when
$S^{-1/2}\le c_B\nu$, without any openness or connectedness
assumption on the supports.

Apply the smaller-scale product estimate in every cube.  Its geometric
coefficient is the product of the three square roots of incident
energies.  After division of coordinates by $S$, the padded tubes
have width $\Delta=S^{-1/2+\theta}$.  Enlarge them by a fixed factor,
so an incident tube contains the whole smaller grid cube.
Equation~\eqref{eq:transverse-kakeya}, divided by that cube's volume
$S^{-3/2}$, gives the bound
\[
 C_{B,\theta,\sigma}\nu^{-M_\sigma}
 S^{3/2}(S^{-1/2+\theta})^{3-\sigma}
 \prod_i\|v_i\|_2
 \le C_{B,\theta,\sigma}\nu^{-M_\sigma}
 S^\beta\prod_i\|v_i\|_2,\qquad
 \beta=3\theta+\sigma/2,
\]
for the sum of these geometric coefficients.
To bound discarded-field errors, normalize each parent input, use
rapid tails for a discarded factor and the fixed-frequency
$L^\infty$ bound for the others, and expand the three-factor product.
The seven error terms together contribute at most
$C_{B,\theta,N}S^{3-N}\prod_i\|v_i\|_2$, for every prescribed $N$.
This is relative to the current \emph{parent} norms; no lower bound
for a retained local norm is required.  We obtain the recurrence
\begin{equation}\label{eq:transverse-square-root}
 A_B(S,\nu)\le C_{B,\theta,\sigma}\nu^{-M_\sigma}S^\beta
       A_{B+1}(C_B\sqrt S,\nu/2)
       +C_{B,\theta,N}S^{3-N},
 \qquad S^{-1/2}\le c_B\nu.
\end{equation}
Spatial translations only modulate the exact inputs.  Taking $N>3$
and setting $\widetilde A_B=1+A_B$ absorbs the additive error in
the same multiplicative recurrence for $\widetilde A_B$.

Fix the target $\epsilon>0$.  Choose a finite $J$ with
$3\cdot2^{-J}<\epsilon/2$, then choose $\theta,\sigma$ with
$2\beta<\epsilon/2$.  The finitely many radii defined by
$S_0=S$, $S_{j+1}=C_{B+j}\sqrt{S_j}$ are comparable to
$S^{2^{-j}}$, with constants depending only on $B,J$.
If $S\ge(C_{B,J}/\nu)^{2^J}$ for a sufficiently large constant,
all $J$ support-enlargement conditions hold.  Iterating
\eqref{eq:transverse-square-root} and using the terminal trivial
bound $\widetilde A_{B+J}(R,\nu')\lesssim_{B,J}R^3$ gives
\[
 \widetilde A_B(S,\nu)
 \le C_{B,\epsilon}\nu^{-M_{B,\epsilon}}
 S^{\beta\sum_{j=0}^{J-1}2^{-j}+3\cdot2^{-J}}
 \le C_{B,\epsilon}\nu^{-M_{B,\epsilon}}S^\epsilon.
\]
In the complementary range, the trivial bound $A_B(S,\nu)\lesssim_B
S^3$ is absorbed by increasing the fixed inverse-transversality
exponent to at least $3\cdot2^J$.  This proves
\eqref{eq:transverse-product} for all $S,\nu$.  On a subpower-transverse
sequence the threshold is eventually satisfied for each fixed $J$,
and every intermediate enlarged support remains subpower transverse.
This is a finite version of the restriction bootstrapping in
\cite[Proposition~2.1]{BCT06}; polynomial dependence is supplied by this argument,
not by the statement of their restriction theorem alone.

Finally, fixed-frequency reproducing kernels give, for $z\in Q$,
\[
 \sup_Q|Ev_i|\lesssim_N
       \int(1+|w|)^{-N}|Ev_i(z+w)|\,dw.
\]
Integrate the product over $Q$ and sum.  For each triple of shifts the
three shifted fields are exact extensions on the same domain, with
independently modulated inputs of unchanged norms and supports.  Apply
\eqref{eq:transverse-product} and integrate the Schwartz weights.
\end{proof}

\begin{lemma}[Reduction to weighted incident tubes]
\label{lem:transverse-tubes}
A countersequence to Proposition~\ref{prop:transverse} yields a set
$P$ in a fixed bounded region and three weighted tube families satisfying
\begin{equation}\label{eq:transverse-point-counts}
 |P|=L^{2+o(1)},\qquad
 |P\cap B(z,r)|\le L^{2+o(1)}r^2
             \quad(L^{-1}\le r\le1).
\end{equation}
The points are centers of distinct cubes of side comparable to $L^{-1}$.
The tubes have width $L^{-1+o(1)}$, their coefficient weights obey
$\sum_Tw_i(T)\lesssim1$, and every cross-family direction triple has
determinant at least $L^{-o(1)}$.  At each $p\in P$ there is a retained
incident subfamily of total weight $e_i^*(p)=L^{-4/3+o(1)}$ such that
its normalized weight law satisfies, for fixed $\kappa>0$, $u_0>0$,
\begin{equation}\label{eq:src1}
 \Prob\bigl(\operatorname{dir}T_i\in B(v,L^{-u})\mid X=p\bigr)
       \le L^{-\kappa u+o(1)}
       \qquad(0<u<u_0,\ u\in\mathbb Q).
\end{equation}
This bound is uniform in $p,v$ for each fixed rational $u$.
The retained subfamily may depend on $p$.
\end{lemma}

\begin{proof}
Sort occupied $L$-boxes by their number of selected unit cubes.
A dyadic occupancy class retains $L^{4-o(1)}$ cubes.  The squared
supremum sum for each field on this class is at least
$L^{4/3-o(1)}$.  By Lemma~\ref{lem:fractal} at radius $L^2$, a class
with occupancy at most $L^{2-\delta}$, for fixed $\delta>0$, instead
has squared sum at most $L^{4/3-\delta/3+o(1)}$.  Thus some retained
class has $L^{2-o(1)}$ cubes in each occupied box.  The upper counting
bound gives occupancy at most $L^{2+o(1)}$.  Divide coordinates by
$L^2$ and take the centers of these boxes.  Counting cubes in an
enlarged ball and dividing by the minimum occupancy proves
\eqref{eq:transverse-point-counts}.  Translation of the original
ambient ball merely modulates the exact data.

Apply Lemma~\ref{lem:packets} at packet width $L$ on the time
interval of length $O(L^2)$.  After division of coordinates by $L^2$,
the resulting tubes have width $L^{-1+o(1)}$.  Give each tube the
squared coefficient of its packet as weight.  On each retained box,
packets reaching a padded enlargement give an exact local input and
an error $O(L^{-100})$.  The padding exponent tends sufficiently
slowly to zero; the rapid-tail estimates permit this by diagonalization.
By \eqref{eq:prelim-subseries-energy}, the squared norm of this
exact local input is bounded by the incident coefficient energy
$e_i(p)=\sum_{T\ni p}w_i(T)$ up to a constant.  On the at least
$L^{2-o(1)}$ selected unit cubes in this box the required squared
supremum sum is $L^{-2/3-o(1)}$.  The weighted fractal estimate at
radius $L$ bounds it above by $L^{2/3+o(1)}e_i(p)$, so
\[
 e_i(p)\ge L^{-4/3-o(1)}.
\]
Lemma~\ref{lem:transverse-product}, in its Kakeya form and with
$\Delta=L^{-1+o(1)}$, yields
\[
 \sum_{p\in P}\prod_{i=1}^3e_i(p)^{1/2}\le L^{o(1)}.
\]
For fixed $\delta>0$, points at which some
$e_i(p)\ge L^{-4/3+\delta}$ therefore number at most
$L^{2-\delta/2+o(1)}$.  Discarding these points by a diagonal choice
leaves \eqref{eq:transverse-point-counts} and makes all three energies
$L^{-4/3+o(1)}$ uniformly.

It remains to remove concentration in direction without losing the
local lower bound.  Choose $0<\kappa<1/3$ and $0<u_0<1/4$.
For each fixed rational $0<u<u_0$, put $q=L^{-u}$ and partition
directions into bounded-overlap bins of diameter comparable to $q$.
Delete incident bins having weight greater than $q^\kappa e_i(p)$.
There are $O(q^{-\kappa})$ such bins.  On the original box of radius
$O(L)$, Lemma~\ref{lem:cap} applies with $S=L$; the packet frequency
enlargement $O(L^{-1})$ is smaller than the bin width $q$.
Cauchy--Schwarz over the deleted bins then bounds the squared
supremum sum of their combined field by
\[
 L^{2/3+o(1)}q^{1/3-\kappa}e_i(p).
\]
Here the chosen unit cubes have local density $L^{o(1)}$, and the
energies of all deleted bins sum to $O(e_i(p))$.  The estimate also
holds for arbitrary subfamilies of these bins.  For every fixed $u>0$
this is negligible relative to $L^{-2/3-o(1)}$.

Delete successively for any fixed finite list of rational $u$'s,
assigning a packet deleted more than once to just one bin.  The
triangle inequality in the square-sum norm shows that the surviving
field retains its lower local norm.  The unrestricted local estimate
then gives surviving weight $e_i^*(p)=L^{-4/3+o(1)}$.  Every surviving
bin has original weight at most $q^\kappa e_i(p)$; normalizing by
$e_i^*(p)$ costs only $L^{o(1)}$.  Boundedly many neighboring bins
cover a ball, and a diagonal choice proves \eqref{eq:src1} at all fixed
rational scales.  All estimates were uniform in $p$, so this
point-dependent pruning supplies legitimate conditional laws without
requiring a single globally pruned family.
\end{proof}

\subsection{Information at spatial scales}

From now on only the abstract configuration of
Lemma~\ref{lem:transverse-tubes} is used.  Draw $X$ uniformly from $P$.
Given $X$, draw six tubes $T_i,T_i'$ independently, two from each
retained family.  A tube label includes its exact direction.  For
rational $0<s\le1$, let $X_s$ be the cell containing $X$ in nested
dyadic grids of side comparable to $L^{-s}$; use the trivial partition
at $s=0$.  For a discrete variable $Y$ and arbitrary conditioning $V$,
write
\[
 I_L(Y\mid V)=-\log_L\Prob(Y=Y(\omega)\mid V)
\]
for its information evaluated at the sampled labels.  Conditional
probabilities may be fixed in any version on null sets.
This is the exponent of the reciprocal probability of the realized label,
rather than its expectation: a sampled label of conditional probability
$L^{-a+o(1)}$ has information $a+o(1)$.  We retain these sampled
exponents in one joint limiting law so that comparisons at different
spatial scales refer to the same configuration.

\begin{lemma}[Limiting information calculus]
\label{lem:transverse-information}
After a subsequence, all the countably many information variables used
below have a common limiting joint law.  For spatial, coordinate, and
direction partitions with at most $L^{C+o(1)}$ possible labels,
their conditional informations are uniformly integrable and have
bounded limiting values.  In this law the chain rule remains exact and
\begin{equation}\label{eq:transverse-conditioning}
 I(Y\mid V)\ge I(Y\mid V,W)\quad\hbox{almost surely}.
\end{equation}
If, given $Z,V$, the possible values of $Y$ on an event $E_L$
number at most $L^{\gamma+o(1)}$, include its indicator in the
joint limit.  On the event where the limiting indicator equals one,
$I(Y\mid Z,V)\le\gamma$.  Alternatively, an event specified by
a finite system of sampled inequalities can be treated by open-set
implications with fixed slack before that slack tends to zero; no
identification of arbitrary boundary events is implicit.  Consequently
variables determining one
another up to $L^{o(1)}$ possibilities have identical limiting
information, also after additional conditioning.  Finite systems of
limiting inequalities holding with positive probability pass to
finite $L$ with any fixed positive slack.
\end{lemma}

\begin{proof}
A variable with at most $L^{C+o(1)}$ labels satisfies, uniformly in
conditioning,
\[
 \Prob\{I_L(Y\mid V)>C+\eta\}\le L^{-\eta+o(1)}.
\]
The same counting argument with a varying threshold proves uniform
integrability.  For the conditioning assertion put
$p=\Prob(Y\mid V)$ and $q=\Prob(Y\mid V,W)$ at the sampled labels.
Conditional on $(V,W)$, the expectation of $p/q$ is at most one.
Thus
\[
 \Prob\{I_L(Y\mid V)-I_L(Y\mid V,W)<-\eta\}
       \le L^{-\eta}.
\]
This proves the almost-sure sign in the limit; no finite-scale
pointwise monotonicity is asserted.  The chain rule is pointwise
already at finite scale.  The determination rule follows by summing
over the stated small set of possible labels.  This argument also
works on restricted events, since it bounds the total mass of the
bad low-probability labels in the event's permitted set.  If the
event indicator is included in the joint limit, apply weak convergence
to the open set where it is one and the information is strictly
above $\gamma+\eta$, and then let $\eta\downarrow0$.  For events
defined through other limiting variables, the same argument uses
strict inequalities and slack, as detailed in the zero-ratio
application below.

Tightness and diagonal weak convergence on a countable product give
the joint law.  Include in this product all rational block and prefix
parameters, all finitely many choices of reference and replacement
tubes, their direction bins, coordinate grids, projections, the
information differences used below, and the angular quantities to be
introduced.  Direction grids use deterministic representatives of
occupied bins.  Realized bases and inverse bases have norms
$L^{o(1)}$, so every coordinate partition has at most $L^{C+o(1)}$
labels.  Standard open-set inequalities for weak convergence justify
the assertion about finite systems after enlarging their bounds by
slack.  Uniform integrability permits passage of expectations in the
cost estimates below.
\end{proof}

\begin{lemma}[The spatial profiles]\label{lem:transverse-profiles}
In the limiting law set
\[
 F(s)=I(X_s),\qquad G_i(s)=I(X_s\mid T_i),\qquad
 G_i'(s)=I(X_s\mid T_i').
\]
These functions extend from rational scales to nondecreasing Lipschitz
functions on $[0,1]$, vanishing at zero.  Their Lipschitz constants
are respectively $3$ and $1$.  Almost surely,
\begin{equation}\label{eq:transverse-endpoint-information}
 F(s)\ge2s,\qquad F(1)=2,\qquad
 G_i(1),G_i'(1)\ge\tfrac23.
\end{equation}
\end{lemma}

\begin{proof}
The exact chain increment of $G_i$ between $s$ and $s+t$ is
$I(X_{s+t}\mid X_s,T_i)$.  Inside a parent, the tube meets at most
$L^{t+o(1)}$ child cubes, including its padding.  The determination
rule bounds this increment between zero and $t$; the unrestricted
spatial count gives the bounds zero and $3t$ for $F$.  Rational
comparisons simultaneously yield the Lipschitz extensions.

The ball count in \eqref{eq:transverse-point-counts} gives every
$X_s$-cell mass at most $L^{-2s+o(1)}$, proving $F(s)\ge2s$.
An $X_1$-cell contains boundedly many points and $X$ is uniform, so
$F(1)=2$.  Finally the retained weights give
\[
 \Prob(X=p,T_i=T)
   =|P|^{-1}\frac{w_i(T)\mathbf1_{\{T\text{ retained at }p\}}}
                      {e_i^*(p)}
   \le L^{-2/3+o(1)}w_i(T).
\]
The probability of sampling a tube whose marginal is smaller than
$L^{-\eta}w_i(T)$ is at most
$L^{-\eta}\sum_Tw_i(T)=O(L^{-\eta})$.  Outside this event, every
conditional point mass is at most $L^{-2/3+\eta+o(1)}$.
Letting $\eta\downarrow0$ proves the conditional endpoint bounds,
also for all primed tubes.  This uses only summability of weights,
so exact tube labels need not have bounded cardinality.
\end{proof}

\subsection{Direction costs and equality of the profiles}

To compare the three tube-conditional profiles, we use coordinates
adapted to one reference tube from each family.  Specifying those
coordinates requires revealing coarse direction bins.  This can lower
the spatial information; the costs below measure that loss, and their
telescoping bound will let the loss tend to zero with the block size.

Fix a rational block $[s,s+b]\subset[0,1]$ and rational $b_*\ge b$.
Let $C=X_s$ and let $D$ record direction bins of precision $L^{-b_*}$
for a fixed sublist of the six tubes, including one reference tube
from each family.  Within the parent subtract its center, multiply
coordinates by $L^s$, and use the approximate basis supplied by the
three reference bins.  Its norms and inverse norms are $L^{o(1)}$.
For rational $0<t\le b$ let $Z_i$ be its $i$th coordinate grid index
at resolution $L^{-t}$, and write $Z=(Z_1,Z_2,Z_3)$ and
$Z_{\widehat i}=(Z_j)_{j\ne i}$.  These coordinate grids are nested
in $t$.  When extra bins are added to $D$, the basis is still selected
using only the three reference bins, so the coordinate variables do
not change.

\begin{lemma}[Block comparison and telescoping costs]
\label{lem:transverse-costs}
Write $J=F(s+t)-F(s)$ and $J_i=G_i(s+t)-G_i(s)$ for the chosen
reference tubes.  Define
\[
 \ell=I(X_{s+t}\mid C)-I(X_{s+t}\mid C,D),\qquad
 \ell_i=I(X_{s+t}\mid C,T_i)-I(X_{s+t}\mid C,T_i,D).
\]
All these limiting costs are nonnegative, and
\begin{equation}\label{eq:src2}
 \sum_{i=1}^3(J_i-\ell_i)
 \le\sum_{i=1}^3I(Z_i\mid Z_{\widehat i},C,D)
 \le I(Z\mid C,D)=J-\ell.
\end{equation}
For disjoint blocks or disjoint prefixes, with the same $D$, the
sum of expectations of each type of cost is $O(b_*)$.  The same
statement holds with any one fixed other tube in place of $T_i$.
All coordinatewise deficits obtained by using any ordering in the
second inequality of \eqref{eq:src2} are nonnegative.
\end{lemma}

\begin{proof}
The spatial child and the three coordinate labels determine one
another up to $L^{o(1)}$ choices.  A reference tube determines its
two off-axis coordinates to $L^{o(1)}$ choices: both its width and
the error in its approximate direction, relative to the parent,
are at most $L^{-t+o(1)}$.  Conditional information calculus therefore
gives
\[
 J_i-\ell_i=I(Z_i\mid C,D,T_i)
   =I(Z_i\mid Z_{\widehat i},C,D,T_i)
   \le I(Z_i\mid Z_{\widehat i},C,D).
\]
For each ordering of the coordinates, its chain-rule term conditions
on a subset of $Z_{\widehat i}$, hence dominates the last expression.
Sum the three terms to obtain \eqref{eq:src2}.

At finite scale, the expected cost is the conditional mutual
information, divided by $\log L$, between the child and $D$ given
the parent (and, when present, the same fixed tube).  Along the nested
spatial filtration these mutual informations telescope.  Inserting
the omitted gaps between disjoint blocks only adds nonnegative
expected terms.  Their sum is at most $H(D)/\log L=O(b_*)$, because
there are $L^{O(b_*)}$ direction-bin lists.  Fixed tube conditioning
is never changed within this telescoping argument.  Passage to the
limit follows from Lemma~\ref{lem:transverse-information}.
\end{proof}

\begin{lemma}[Additivity and replacement invariance]
\label{lem:transverse-additivity}
Almost surely, simultaneously for every choice of one of the two
tubes in each family,
\[
 F=G_1+G_2+G_3,\qquad G_i'=G_i,\qquad G_i(1)=\tfrac23.
\]
In any block above, including a prefix and with any extra direction
bins in $D$,
\begin{equation}\label{eq:src3}
 \begin{split}
 I(Z_i\mid C,D)&=J_i+O\left(\sum_k\ell_k\right),\\
 I(Z_i,Z_j\mid C,D)&=J_i+J_j+
                       O\left(\sum_k\ell_k\right).
 \end{split}
\end{equation}
The absolute constants in these error bounds are fixed.
\end{lemma}

\begin{proof}
Partition a fixed rational spatial interval into blocks of length
at most $b=b_*$.  By \eqref{eq:src2}, the positive part of the excess
of the summed $G_i$ increment over the $F$ increment has expectation
$O(b)$.  Let $b\downarrow0$ to obtain the interval inequality almost
surely.  At $[0,1]$, \eqref{eq:transverse-endpoint-information}
forces equality and $G_i(1)=2/3$.  The interval deficits are
nonnegative and additive, so their zero full-interval total forces
all rational subinterval deficits to vanish.  Continuity gives
$F=\sum_iG_i$.  Repeat for the finitely many replacement triples;
subtracting identities gives $G_i'=G_i$.

For any coordinate ordering, its $i$th chain term dominates
$I(Z_i\mid Z_{\widehat i},C,D)$ and then $J_i-\ell_i$.
The sum of the surpluses is exactly
$\sum_i\ell_i-\ell\le\sum_i\ell_i$, using $J=\sum_iJ_i$.
Take a chosen coordinate, or chosen pair, first in the ordering.
This gives both bounds in \eqref{eq:src3}, with each individual
error controlled by the total cost.
\end{proof}

\begin{lemma}[One common profile]\label{lem:transverse-common-profile}
Almost surely,
\begin{equation}\label{eq:src4}
 G_1=G_2=G_3=:G,\qquad F=3G.
\end{equation}
In particular $0\le\dot G\le1$, $G(1)=2/3$, and $3G(s)\ge2s$.
\end{lemma}

\begin{proof}
In the exact basis of $T_1,T_2,T_3$, normalize the $i$th coordinate
of the direction of $T_i'$ to one, and call its other two ratios
$a_j$ ($j\ne i$).  Replacement transversality and Cramer's rule
bound the normalizing denominator, its inverse, and all resulting
ratios by subpowers.  Include in the joint limit the values
$\min\{1,\max\{0,-\log_L|a_j|\}\}$, with value one at zero.
The two limiting values cannot both be positive: that would put
$T_i'$ within a fixed positive-power angular ball about $T_i$,
with probability tending to zero by \eqref{eq:src1} and conditional
independence given $X$.

Let $r_j$ be the limiting clipped value for $a_j$.  We make the
restricted-event comparison on $\{r_j=0\}$ precise.  Fix a rational
$0<\eta<t/2$.  At finite scale, $r_{j,L}<\eta/2$ implies
$|a_j|>L^{-\eta/2}$.  The approximate basis fixed by $D$ differs
from the exact basis by $L^{-b_*+o(1)}$, including the subpower
condition numbers.  Its corresponding coefficient thus has
magnitude at least $L^{-\eta}$ for large $L$.  Given
$(T_i',Z_j,C,D)$ alone, possible child positions on this event
number at most $L^{O(\eta)+o(1)}$: the approximate $j$th coordinate
parametrizes the tube with inverse slope at most $L^{\eta+o(1)}$.
This permitted set depends only on those displayed labels.  Exact
reference directions need not be added to the conditioning.

Apply the restricted determination bound and the information
conditioning rule.  The probability that $r_{j,L}<\eta/2$ and
\[
 I_L(X_{s+t}\mid C,D,T_i')-I_L(Z_j\mid C,D)>C\eta+\eta
\]
tends to zero, with a fixed absolute $C$.  The joint limit therefore
gives the corresponding comparison on the open event
$\{r_j<\eta/2\}$ with error at most $(C+1)\eta$; one can first
use any strict upper error bound and then decrease it.  This is
an implication about the support of the joint limiting law and
requires no coupling of finite-scale and limiting events.  Let
rational $\eta\downarrow0$.  On $\{r_j=0\}$ we obtain
\[
 G_i'(s+t)-G_i'(s)-\ell_i'
       \le I(Z_j\mid C,D),
\]
where $D$ includes the alternative bin and $\ell_i'$ is its analogous
cost.  Replacement invariance and \eqref{eq:src3} now compare the
$i$ and $j$ profile increments up to the costs.

Take dyadic partitions of the exponent interval, with block size
$b=b_*$.  On each block distribute its cost uniformly over its
length.  The integral of this density is the cost; hence its expected
integral is $O(b)$ by Lemma~\ref{lem:transverse-costs}.  Summability
over dyadic $b$ shows that all normalized costs tend to zero for
almost every pair consisting of a path and an exponent.  At
simultaneous differentiability points it follows that
$\dot G_i\le\dot G_j$ whenever the corresponding limiting ratio
information is zero.  The restricted-event determination just used
is applied before differentiation, so no uniform positive lower
bound for the coefficient is required.

For almost every path each vertex $i$ therefore has an edge $i\to j$
for which $\dot G_i\le\dot G_j$ almost everywhere.  Since both
integrals are $2/3$, this forces $G_i=G_j$.  Every equivalence class
of equal profiles has at least two members.  With three families
there is only one such class, proving \eqref{eq:src4}.
\end{proof}

\subsection{Projection growth on a block}

The next input is the robust form of the discretized projection theorem:
one direction improves the projection of every sufficiently large subset.
That quantifier is essential because the subset we construct will depend
on the direction.

\begin{lemma}[Robust planar discretized projection]\label{lem:transverse-projection}
Fix $0<d<1$ and $\rho_0>0$.  There are $\varepsilon_{\rm pr}>0$
and $\eta_{\rm pr}>0$ with the following property for all sufficiently
small $\Delta$.  Suppose $A\subset B(0,1)\subset\R^2$, with
$\mathcal N_\Delta$ denoting covering number by radius-$\Delta$ balls,
satisfies
\[
 \Delta^{-2d+\eta_{\rm pr}}\le\mathcal N_\Delta(A)
       \le\Delta^{-2d-\eta_{\rm pr}},\qquad
 \mathcal N_\Delta(A\cap B(z,r))
       \le\Delta^{-\eta_{\rm pr}}r^{\rho_0}
                                  \mathcal N_\Delta(A)
\]
for $\Delta\le r\le1$.  Let $\sigma$ be a probability on projection
lines obeying
$\sigma(B(V,r))\le\Delta^{-\eta_{\rm pr}}r^{\rho_0}$ on the same
range.  Some $V\in\supp\sigma$ then satisfies
\[
 \mathcal N_\Delta(\pi_V A')\ge\Delta^{-d-\varepsilon_{\rm pr}}
 \quad\hbox{for every }A'\subset A\hbox{ with }
 \mathcal N_\Delta(A')\ge\Delta^{\eta_{\rm pr}}
                                      \mathcal N_\Delta(A).
\]
One may decrease both the tolerance and the gain.
\end{lemma}

\begin{proof}
This is the $n=2$, $m=1$, $\alpha=2d$ case of
\cite[Theorem~1]{He20}, a general form of the planar Bourgain
discretized projection theorem.  In dimension two its Grassmannian
nonconcentration condition is the displayed angular-ball condition.
That theorem gives the conclusion simultaneously for every large
subset, for a direction set of probability at least one minus a
positive power of $\Delta$.  Decreasing the tolerance relative to
its exponent gives the stated formulation.
\end{proof}

Take $T_1,T_2,T_3$ as reference tubes and write the direction of $T_1'$
in their exact basis as $(1,a_2,a_3)$.  Given $X$ and these three
exact references, redraw $T_1'$ from its original conditional law.
Let $h_j(u)$, for rational $0<u<u_0$, be the limiting value of the
negative base-$L$ logarithm of the resulting marginal mass in the
ball of radius $L^{-u}$ centered at the sampled $a_j$.

\begin{lemma}[Angular marginal information]\label{lem:transverse-angular-information}
The functions $h_2,h_3$ are nonnegative and nondecreasing on their
rational domains, and almost surely
\begin{equation}\label{eq:transverse-angular-sum}
 h_2(u)+h_3(u)\ge\kappa u\qquad(0<u<u_0,\ u\in\mathbb Q).
\end{equation}
Their values agree with those of marginal coordinate-cell information
at comparable resolution; the analogous joint assertion also holds.
\end{lemma}

\begin{proof}
The exact coordinate map and its inverse have subpower norm on the
realized support.  A joint ball in $(a_2,a_3)$ therefore pulls back
to at most two direction balls of radius $L^{-u+o(1)}$.  The
alternative direction law, conditional on $X$ and the exact
references, is still its original law given $X$.  Equation
\eqref{eq:src1}, enlarged by a subpower using covering, bounds this
joint ball mass by $L^{-\kappa u+o(1)}$.

Here ball and grid informations have the same limit.  Conditional
on the exact point and references, take grids at comparable sizes.
If $p_z$ is the mass of a sampled cell and $q_z$ the total mass of
its boundedly many neighbors, then
\[
 \sum_{z:q_z>L^\eta p_z}p_z
       \le L^{-\eta}\sum_zq_z\lesssim L^{-\eta}.
\]
A boundedly finer grid and the determination rule handle balls not
contained in one grid cell.  The argument is identical for joint
and marginal grids.  The chain rule and
\eqref{eq:transverse-conditioning} bound joint cell information by
the sum of marginal informations.  This proves
\eqref{eq:transverse-angular-sum}; monotonicity follows directly
from nesting of centered balls.
\end{proof}

\begin{lemma}[The projection block test]\label{lem:transverse-block-test}
Fix $0<d<1$, $d_0>0$, and $\kappa_0>0$.  There is $\zeta>0$,
depending only on these parameters, for which the following
conditions cannot hold on an event of positive limiting probability.
Fix rational $[s,s+b]\subset[0,1]$, with $b<u_0$, rational
$b_*\ge b$, and $j\in\{2,3\}$.  At direction precision $L^{-b_*}$
let $D_0$ be the reference triple of bins and let $D_1$ add the
alternative bin.  Require
\begin{enumerate}[label=\textnormal{(\roman*)}]
 \item $|G(s+b)-G(s)-db|\le\zeta b$, and
       $G(s+u)-G(s)\ge d_0u-\zeta b$;
 \item all costs in Lemma~\ref{lem:transverse-costs}, for both
       $D_0$ and $D_1$ and each reference tube or $T_1'$, are at most
       $\zeta b$;
 \item $h_j(u)\ge\kappa_0u-\zeta b$.
\end{enumerate}
The prefix, cost, and angular requirements need be imposed only on
a fixed finite rational grid of $0<u\le b$ containing $b$, with
mesh at most a sufficiently small fixed multiple of $\zeta b$.
Decreasing $\zeta$ preserves the conclusion.  All power-error
constants are independent of $s,b,b_*$; the meshes can be fixed
in relative coordinates $u/b$, and their finite cardinalities
only affect fixed constants in the later cost sums.
\end{lemma}

\begin{proof}
Put $K=L^b$ and use the same reference coordinates for $D_0,D_1$.
Let $Y=(Z_1,Z_j)$ at the endpoint or a tested prefix.  Equations
\eqref{eq:src3}--\eqref{eq:src4} imply, with either conditioning
$C,D_0$ or $C,D_1$,
\begin{equation}\label{eq:transverse-pair-information}
 I(Y_b\mid C,D_a)=2db+O(\zeta b),\qquad
 I(Y_u\mid C,D_a)\ge2d_0u-O(\zeta b),\quad a=0,1.
\end{equation}
Let $\widetilde a_j$ be the approximate direction ratio fixed by
$D_1$, and let $\Pi(Y_b)$ record the interval of length $K^{-1}$
containing $z_j-\widetilde a_jz_1$ at the pair-cell center.
The exact and approximate ratios differ by at most
$L^{-b_*+o(1)}$.  Given $T_1',C,D_1$, this projection ranges over
$L^{o(1)}$ cells.  Indeed, the approximate basis is already fixed
by $D_0$, and the exact central line of $T_1'$ is part of its
label.  Its slope differs from the representative slope by
$L^{-b_*+o(1)}$, while its width in parent coordinates is at most
$L^{s-1+o(1)}\le L^{-b+o(1)}$.  Along a segment of length
$L^{o(1)}$ both errors cause only $L^{-b+o(1)}$ projection
variation.  Hidden exact reference directions do not enter this
coordinate system.  Also $Y_b$ and $X_{s+b}$ determine each other
to $L^{o(1)}$ choices with these labels: coordinate one parametrizes
the alternative tube with subpower distortion.  Consequently
\begin{align*}
 I(Y_b\mid\Pi,C,D_1)
 &\ge I(Y_b\mid\Pi,C,D_1,T_1')\\
 &=I(X_{s+b}\mid C,D_1,T_1')
   \ge G(s+b)-G(s)-\zeta b.
\end{align*}
Since $\Pi$ is determined by $Y_b,D_1$, its chain rule together
with \eqref{eq:transverse-pair-information} gives
\begin{equation}\label{eq:transverse-projection-information}
 I(\Pi\mid C,D_1)\le db+O(\zeta b).
\end{equation}
All projection information variables are among those included in
the joint limit.

We detail the passage from these sampled inequalities to one common
set and one direction measure.  Enlarge the inequalities by fixed
slack and pass to finite $L$.  There is a success event of probability
at least $p_*>0$ along a subsequence.  Fix parent and reference-bin
labels $(c,d_0^{\rm bin})$ whose conditional success probability
$p$ is at least $p_*$.  No lower bound on the probability of this
individual fiber is needed.  Denote its conditional law by $\mu$.
Absorb all fixed slack into $\alpha=C\zeta+o(1)$, enlarging the
absolute constant $C$ as necessary.  On success,
\begin{equation}\label{eq:transverse-fiber-masses}
 \begin{split}
 K^{-2d-\alpha}\le\mu(Y_b=y)&\le K^{-2d+\alpha},\\
 K^{-2d-\alpha}\le\mu(Y_b=y\mid A_{\mathrm{bin}}=a)
                                     &\le K^{-2d+\alpha},\\
 \mu(Y_u=y_u)&\le K^{\alpha}L^{-2d_0u},\\
 \mu(\Pi=\pi\mid A_{\mathrm{bin}}=a)&\ge K^{-d-\alpha},
 \end{split}
\end{equation}
where $A_{\mathrm{bin}}$ is the additional direction bin.  These
are statements about the indicated labels whenever they have any
successful realization.

Let $A$ consist of the centers of all endpoint pair cells with a
successful realization in this fixed parent/reference fiber.
The first line of \eqref{eq:transverse-fiber-masses} and total
success mass at least $p$ give
\begin{equation}\label{eq:transverse-common-set-size}
 pK^{2d-\alpha}\le |A|\le K^{2d+\alpha}.
\end{equation}
A prefix cell containing a successful endpoint has mass at most
the third line of \eqref{eq:transverse-fiber-masses}.  The nested
grids put all mass of that endpoint cell in its prefix.  Dividing
by the minimum endpoint mass bounds the number of good endpoints
in that prefix.  A ball of radius $r$ meets boundedly many cells
at a comparable coarser tested prefix size; the finite mesh costs
$K^{O(\zeta)}$.  Thus, for $K^{-1}\le r\le1$,
\begin{equation}\label{eq:transverse-common-set-counts}
 \frac{\mathcal N_{K^{-1}}(A\cap B(z,r))}
      {\mathcal N_{K^{-1}}(A)}
       \lesssim p^{-1}K^{O(\zeta)+o(1)}r^{d_0}.
\end{equation}
In fact the preceding argument gives $r^{2d_0}$ before weakening
the exponent and absorbing rounding.  For the untested shortest
prefixes the allowed power factor suffices.

Retain bins $a$ whose conditional success probability is at least
$p/2$.  The successful mass in these bins is at least $p/2$.
For each retained bin, let $A_a\subset A$ be its successful endpoint
cells.  The conditional upper pair-cell mass gives
\[
 |A_a|\ge(p/2)K^{2d-\alpha},\qquad
 |A_a|/|A|\ge(p/2)K^{-2\alpha}.
\]
The lower projection-cell mass in
\eqref{eq:transverse-fiber-masses} gives
\begin{equation}\label{eq:transverse-small-projections}
 \mathcal N_{K^{-1}}\bigl(\pi_{\widetilde a_j(a)}A_a\bigr)
          \le K^{d+O(\zeta)+o(1)},
\end{equation}
where initially $\pi_{\widetilde a_j}(z_1,z_j)
=z_j-\widetilde a_jz_1$.  Passing to a unit vector for this
projection only shrinks interval lengths.

It remains to put a nonconcentrating probability on these directions.
Restrict the actual samples to success and the retained bins, normalize,
and push forward to the line spanned by $(-\widetilde a_j,1)$.
This costs at most $2/p$ in mass.  To estimate a slope ball, first
condition on the exact point $X$ and exact reference tubes, before
any restriction by the additional bin.  These labels determine
$(C,D_0)$, so fixing that fiber does not change the original
alternative law once these exact labels are given.  If that ball contains any
successful alternative, the centered marginal bound at this
alternative, supplied by (iii), bounds the \emph{original}
alternative probability of the entire ball after a fixed radius
enlargement.  Choosing a coarser tested radius costs
$K^{O(\zeta)+o(1)}$.  Therefore its restricted mass is at most
\[
 K^{O(\zeta)+o(1)}r^{\kappa_0}\qquad(K^{-1}\le r\le1).
\]
Integrate this bound over the exact points and references under the
fixed fiber, and then normalize.  This operation introduces no
reciprocal probability of an individual additional bin; conditional
independence supplies the original alternative law at the stage
where the angular estimate is used.

All slopes have magnitude $L^{o(1)}$.  The sine-angle formula shows
that the slopes in a projective ball of radius $r$ occupy an interval
of length $rL^{o(1)}$; exact-to-approximate slope errors are
$K^{-1}L^{o(1)}$.  Covering and radius rounding thus give, for the
pushforward probability $\sigma$,
\begin{equation}\label{eq:transverse-direction-measure}
 \sigma(B(V,r))\le K^{O(\zeta)+o(1)}r^{\kappa_0},
                   \qquad K^{-1}\le r\le1.
\end{equation}
The coordinate range of $A$ is at most $L^{o(1)}$.  Rescale it into
the unit ball.  At resolution $K^{-1}$ this merges or covers at most
$L^{o(1)}$ grid cells per ball and changes every displayed power
error by $o(1)$ only.  Use covering numbers after this rescaling.

Choose $0<\rho_0<\min\{d_0,\kappa_0\}$, apply
Lemma~\ref{lem:transverse-projection}, and then take $\zeta$ so
small that every $O(\zeta)$ error is below its tolerance and its
projection gain.  Equations
\eqref{eq:transverse-common-set-size},
\eqref{eq:transverse-common-set-counts}, and
\eqref{eq:transverse-direction-measure} meet its hypotheses.
For every direction in the finite support of $\sigma$, however,
\eqref{eq:transverse-small-projections} gives a subset of the
\emph{same} $A$ of relative covering number
$K^{-O(\zeta)-o(1)}$ with too small a projection.  The robust
large-subset conclusion gives the contradiction.
\end{proof}

\subsection{Selecting angular scales and differentiating}

The sum bound \eqref{eq:transverse-angular-sum} does not say that one
fixed marginal is nonconcentrated at every radius.  The following
finite choice of comparable scales supplies exactly what the block
test needs.

\begin{lemma}[Finite angular alternatives]\label{lem:transverse-angular-choice}
Fix $0<d<1$ and $d_0>0$.  There is a fixed finite list of positive
rational ratios, with minimum $c_*>0$, such that for every sufficiently
small rational $b_0$, almost every path has a choice $b'$ from this
list times $b_0$ and a marginal $j\in\{2,3\}$ satisfying the angular
hypothesis of Lemma~\ref{lem:transverse-block-test}, with spare slack,
for block lengths in $[b'/4,b']$.  Each alternative has its own fixed
positive $\kappa_0$, tolerance, and finite relative mesh.  All choices
are independent of $b_0$, and remain valid under sufficiently small
variations of the spatial increment parameters.
\end{lemma}

\begin{proof}
First obtain a sufficiently reduced block-test tolerance $\zeta_2$
for $\kappa_0=\kappa/4$.  Choose a positive rational
$c_1\ll\kappa\zeta_2$, and then a sufficiently reduced tolerance
$\zeta_1$ for $\kappa_0=c_1/2$.
Test $h_2(u)\ge c_1u$ on a sufficiently fine finite rational grid
from $c\zeta_1b_0$ to $b_0$, where $c>0$ is fixed and small.
If all tests pass, monotonicity fills the mesh gaps; below the
first grid point nonnegativity gives the required inequality with
the allowed additive error.  Take $b'=b_0$ and $j=2$.

Otherwise some grid point $v$ satisfies $h_2(v)<c_1v$.
For rational $u\le v$, monotonicity and
\eqref{eq:transverse-angular-sum} yield
\[
 h_3(u)\ge\kappa u-h_2(u)>\kappa u-c_1v
            \ge\frac\kappa2u
       \quad\left(\frac{2c_1v}{\kappa}\le u\le v\right).
\]
Take $b'=v$, $j=3$.  Because $c_1\ll\kappa\zeta_2$, the missing
smaller radii meet the weaker $\kappa/4$ bound with additive error
$\zeta_2b$ for every $b\in[v/4,v]$.  Small enough mesh and cutoff
constants leave spare slack in both cases.  The possible $v/b_0$
form a fixed finite list bounded away from zero.  Reduce all
constants once more to permit the stated small spatial variations.
\end{proof}

\begin{lemma}[Binary derivative]\label{lem:transverse-binary}
For almost every path, $\dot G(s)\in\{0,1\}$ for almost every
$s\in[0,1]$.
\end{lemma}

\begin{proof}
Suppose otherwise.  Choose $0<d<1$ and $0<d_0<d$ such that a
sufficiently small neighborhood of $d$ contains $\dot G(s)$ on a
set of positive path-times-exponent measure.  This is possible by
choosing an interior point of the essential support.  Apply
Lemma~\ref{lem:transverse-angular-choice} for these parameters.
For dyadic $b_0\downarrow0$, and each of its finite alternatives,
form a rational mesh of test blocks of comparable length, so that
each interior exponent is contained in an appropriate block, and
include each block's finite prefix grid.  Use direction precision
$b_*=b_0$ for all alternatives.  Within each size and prefix type,
the candidate intervals have bounded overlap, with constants
independent of $b_0$.  Explicitly, choose endpoints on one lattice
of spacing $q b_0$, where the positive rational $q$ is fixed below
all the relative mesh tolerances, and allow all candidate lengths
between $c_*b_0/4$ and $b_0$.  A point belongs to at most
$O(q^{-2})$ such candidates.  Every tested prefix is contained in
its candidate, and the number of prefix types is fixed.  Color each
of these interval families into finitely many disjoint lists.

Let $\mathcal C_I$ be the sum of all reference, replacement, and
extra-bin costs used by candidate block $I$.  By
Lemma~\ref{lem:transverse-costs},
\begin{equation}\label{eq:transverse-mesh-cost}
 \sum_I\E\mathcal C_I\lesssim b_0.
\end{equation}
All candidate lengths lie between fixed positive multiples of $b_0$.
For the nonnegative function
\[
 R_{b_0}(\omega,s)=
       \sum_{I\ni s}\mathcal C_I(\omega)/b_0
\]
we have $\E\int_0^1R_{b_0}(\omega,s)\,ds\lesssim b_0$:
integration cancels the normalizing block length.  Dyadic summability
therefore gives $R_{b_0}\to0$ almost everywhere.  In particular,
all tested costs in a block containing such an $s$ are $o(b_0)$.

At a differentiability point $s$, the approximation
$G(y)=G(s)+\dot G(s)(y-s)+o(b_0)$ is uniform for
$|y-s|\le Cb_0$.  Hence all increments and finitely many prefix
increments in these nearby blocks are approximated to $o(b_0)$.
If $\dot G(s)$ lies in the chosen sufficiently small neighborhood
of $d$, the spatial hypotheses of the block test hold eventually,
with its reduced tolerance.  An angular alternative supplies its
angular hypotheses, and the preceding cost estimate supplies the
remaining ones.  Only countably many fixed rational blocks and
finite tests are involved.  Positive path-times-exponent measure
therefore implies that one fixed test holds on a positive-probability
event, contradicting Lemma~\ref{lem:transverse-block-test}.
A countable cover of interior derivative values proves the claim.
\end{proof}

\subsection{Hills and finitely many jumps}

Write $E=\{s:\dot G(s)=1\}$ for a path satisfying
Lemma~\ref{lem:transverse-binary}.  Equations
\eqref{eq:transverse-endpoint-information} and \eqref{eq:src4} give
\begin{equation}\label{eq:transverse-full-set}
 |E|=\tfrac23,\qquad 3|E\cap[0,s]|\ge2s\quad(0\le s\le1).
\end{equation}
We next exclude infinitely many changes from a full scale interval
to an empty one.  The block test applies to the balanced intervals
created at such changes, even though their derivatives themselves
are binary.

\begin{lemma}[Balanced subintervals inside a hill]
\label{lem:transverse-hills}
Let $G$ be Lipschitz with derivative $\mathbf1_E$, and let $a<b$
be respectively a density point of $E$ and of its complement.
There is a positive number $b_{\mathrm{hill}}$ such that, for every
$0<r<b_{\mathrm{hill}}$, some $[x,y]\subset(a,b)$ has
\[
 r\le y-x\le2r,\qquad
 G(y)-G(x)=\tfrac12(y-x),\qquad
 G(x+u)-G(x)\ge\tfrac12u\quad(0\le u\le y-x).
\]
For any fixed finite collection of disjoint ordered density pairs,
these intervals can be chosen disjointly at every sufficiently small
common $r$.
\end{lemma}

\begin{proof}
Put $H(s)=G(s)-s/2$.  The density conditions give positive right
slope at $a$ and a positive drop towards $b$ from the left.  Thus
$\max_{[a,b]}H$ exceeds both $H(a)$ and $H(b)$: if the larger
endpoint is $a$ use its right increase, and if it is $b$ use its
left increase.  A component of a suitable strict superlevel set
therefore gives a closed hill $[a',b']\subset(a,b)$ with equal
endpoint heights and strictly larger values inside.

Fix $r<b'-a'$.  Among levels having a closed subinterval of this
hill of length at least $r$, with both endpoints at the level and
values no smaller inside, take the highest.  Such a level exists:
the admissible triples of level and endpoints form a nonempty compact
set.  At the maximizing level $h$ take an admissible interval
$[c,d]$.  No component on which $H>h$ inside it has length greater
than $r$, since such a component contains an interval of length
at least $r$ closed off at a slightly higher level.  The closed
set $\{H=h\}\cap[c,d]$ thus contains both endpoints and has no
gap longer than $r$.  Starting at $c$, take its first point at or
after $c+r$.  Its distance from $c$ is between $r$ and $2r$.
Between these two points $H\ge h$, which gives the asserted total
and prefix identities for $G$.  Apply this construction separately
inside finitely many disjoint hills to obtain the final assertion.
\end{proof}

\begin{lemma}[Finite alternation and a full-to-empty jump]
\label{lem:transverse-finite-jumps}
Almost surely $E$ agrees modulo null sets with a finite union of
intervals and has an interior transition from an interval of full
measure to an interval of zero measure.  For every fixed $\eta>0$
there is a rational block $[s,s+b]$, $0<b<u_0$, on an event of
positive probability, such that
\begin{equation}\label{eq:transverse-jump-block}
 G(s+b)-G(s)=\tfrac b2+O(\eta b),\qquad
 G(s+b/2)-G(s)=\tfrac b2+O(\eta b).
\end{equation}
The absolute error constants can be fixed independently of $\eta$.
\end{lemma}

\begin{proof}
Suppose a path has arbitrarily many disjoint ordered density pairs
of the type in Lemma~\ref{lem:transverse-hills}.  Use the angular
alternatives for $d=1/2$, $d_0=1/4$ at dyadic $b_0\downarrow0$.
For its successful angular size $b'\asymp b_0$, choose the balanced
intervals with lengths between $b'/3$ and $2b'/3$.  Round their
endpoints to a sufficiently fine fixed relative rational mesh.
Lipschitz continuity preserves the endpoint and finite prefix
requirements of the block test with spare slack.  For any prescribed
number of disjoint density pairs this produces that many distinct
spatially and angularly successful candidates at every sufficiently
small $b_0$; the rounding preserves distinctness inside their
separate hills.

Consider all mesh candidates for the fixed finite list of angular
alternatives, and all their prefix and direction costs, with
$b_*=b_0$.  Equation \eqref{eq:transverse-mesh-cost} still holds.
A failure of a candidate's cost test requires one nonnegative cost
at least a fixed positive multiple of $b_0$.  Hence the expected
number $N_{b_0}$ of such failures is bounded uniformly in $b_0$.
Every spatially and angularly successful candidate must fail a cost
test almost surely, by Lemma~\ref{lem:transverse-block-test} and
the countability of the mesh collection.  On the paths just
considered, $N_{b_0}\to\infty$.  Fatou's lemma contradicts their
having positive probability, because $\liminf\E N_{b_0}<\infty$.

Thus the number of alternations among density points of $E$ and
of its complement is bounded on almost every path.  To see the
finite-interval conclusion directly, assign to each ordered density
point the maximum length of an alternating sequence ending there.
These integers are bounded, nondecreasing along the ordered density
points, and strictly increase when the type changes.  Each integer
level therefore has one type and occupies an ordered interval.
The Lebesgue density theorem gives the asserted finite union modulo
null sets.

The prefix inequality in \eqref{eq:transverse-full-set} rules out
an initial interval of zero measure, while $|E|=2/3<1$ forces a
subsequent nonempty interval of zero measure.  There is therefore
an interior full-to-empty jump.  Around that jump choose a short
block whose midpoint differs from the jump by at most $\eta b$,
with its left and right halves inside the adjacent full and empty
intervals except for that discrepancy.  Rational endpoints can be
chosen with $b<u_0$, giving \eqref{eq:transverse-jump-block}.
The countable family of rational blocks covers almost all paths;
some fixed block has positive success probability.  No direction
cost is needed in this choice.
\end{proof}

\subsection{The partitioning contradiction}

A full-to-empty information jump says that almost all spatial branching
occurs in the first half of a block.  At its endpoint, each midpoint
cell contains very few eligible descendants.  We now combine this
with the angular law by drawing two endpoints on the same tube.

\begin{lemma}[No full-to-empty block]\label{lem:transverse-no-jump}
For sufficiently small fixed $\eta>0$, the positive-probability
configuration in \eqref{eq:transverse-jump-block} is impossible.
\end{lemma}

\begin{proof}
Fix such a rational block and put $K=L^b$.  By
\eqref{eq:src4} and the chain rule, its endpoint and midpoint
information increments for $X$ are both $3b/2+O(\eta b)$, while
the conditional endpoint increment given $T=T_1$ is
$b/2+O(\eta b)$.  Pass to finite $L$ with fixed slack, and fix a
parent $C=X_s=c$ with conditional success probability bounded below
by a constant $p_*>0$ along the sequence.  Write
\[
 p_c(z)=\Prob(X_{s+b}=z\mid C=c),\qquad
 q_T(z)=\Prob(X_{s+b}=z\mid T,C=c).
\]
Call an incidence $(T,z)$ good if it has at least one realization
in the success event.  Let $\mathcal Z$ be all endpoint labels with
a good incidence, represented by their centers in rescaled parent
coordinates.  Every such endpoint satisfies
\begin{equation}\label{eq:transverse-good-masses}
 K^{-3/2-O(\eta)}\le p_c(z)\le K^{-3/2+O(\eta)},
 \qquad q_T(z)\le K^{-1/2+O(\eta)}
       \quad\hbox{for a good }(T,z).
\end{equation}
Its midpoint cell has marginal mass at most
$K^{-3/2+O(\eta)}$.  These estimates depend only on the displayed
labels, so they remain true after defining goodness by existence
of a successful realization.  In particular
\begin{equation}\label{eq:transverse-good-counts}
 |\mathcal Z|\le K^{3/2+O(\eta)},\qquad
 \#(\mathcal Z\cap B(z,CK^{-1/2}))\le K^{O(\eta)}
\end{equation}
for every fixed $C$.  For the second bound divide the midpoint
mass upper bound by the endpoint mass lower bound, then cover a
ball by boundedly many midpoint cells.  The central line of a tube
with a good incidence passes within
$\delta=K^{-1}L^{o(1)}$ of the endpoint center.  The parent has a
fixed bounded diameter in these coordinates.

Choose a small fixed $\varepsilon>0$, and set
$P=K^{1/2-\varepsilon}$.  Polynomial partitioning in $\R^3$
\cite[Theorem~4.1]{GK15} gives a nonzero polynomial of degree
$D\lesssim P$ whose complementary open cells each contain at most
\begin{equation}\label{eq:transverse-partition-occupancy}
 O(|\mathcal Z|/P^3)=K^{3\varepsilon+O(\eta)}
\end{equation}
of the centers.  Take a wall of thickness a sufficiently large
constant times $\delta$ about its zero set.  At midpoint resolution
$r=K^{-1/2}$, the number of boxes meeting the wall is
\begin{equation}\label{eq:transverse-wall-count}
 O(Dr^{-2}+D^2r^{-1}+D^3)=O(PK).
\end{equation}
Indeed $\delta=o(r)$.  Intersect the zero set with a fixed ball
strictly larger than the bounded parent; every wall point relevant
to that parent is within the required distance of this restricted
zero set.  Wongkew's hypersurface neighborhood-volume bound
\cite[Main Theorem, p.~178]{W93}, thickened to scale comparable to $r$ in that fixed
enlarged ball, is
$O(Dr+D^2r^2+D^3r^3)$.  Divide by the disjoint box volume $r^3$.
The degree condition $D\lesssim K^{1/2-\varepsilon}$ makes all
three terms in \eqref{eq:transverse-wall-count} at most $O(PK)$.

By \eqref{eq:transverse-good-masses}--\eqref{eq:transverse-good-counts},
the marginal mass of good endpoints lost to the wall is at most
\begin{equation}\label{eq:transverse-wall-mass}
 O(PK)\,K^{O(\eta)}K^{-3/2+O(\eta)}
       =K^{-\varepsilon+O(\eta)}.
\end{equation}
Take $\eta$ small compared with $\varepsilon$, so this tends to
zero.  The enlarged wall ensures that the nearby central-line point
of each surviving incidence lies in the same open cell as its
endpoint center.  A line outside the polynomial zero set crosses
it at most $D$ times, and hence visits at most $D+1=O(P)$ open
cells.  A line contained in the zero set has no off-wall incidence.
Thus every tube's good off-wall endpoints lie in at most $O(P)$
open cells.

Draw $T$ with its original marginal given $C=c$, then draw $z,z'$
independently with law $q_T$.  Let $g(T)$ be the $q_T$ mass of
good off-wall endpoints.  Its mean is at least $p_*/2$ for large
$K$.  Conditional on $T$, Cauchy--Schwarz over its at most $O(P)$
open cells gives probability at least $g(T)^2/O(P)$ that both
draws are good off-wall and belong to the same open cell.
Jensen's inequality therefore gives the unconditional lower bound
\begin{equation}\label{eq:transverse-pair-lower}
 \Prob\{z,z'\text{ good off-wall in the same cell}\mid C=c\}
       \gtrsim P^{-1}=K^{-1/2+\varepsilon}.
\end{equation}
We bound this probability from above by separating near and far
pairs.

For $|z-z'|\lesssim K^{-1/2}$, the second bound in
\eqref{eq:transverse-good-counts} gives at most $K^{O(\eta)}$
possible good second endpoints.  Each costs
$K^{-1/2+O(\eta)}$ under $q_T$, so near pairs contribute at most
\begin{equation}\label{eq:transverse-near-pairs}
 K^{-1/2+O(\eta)}.
\end{equation}
For a farther pair, a tube meeting both endpoint centers to error
$\delta$ must have direction within
$O(\delta K^{1/2})=K^{-1/2+o(1)}$ of their joining line, allowing
two signs.

The conditioning at this point is essential.  Under the double
sampling,
\[
 \Prob(T,z,z'\mid C=c)=\Prob(T\mid C=c)q_T(z)q_T(z').
\]
Thus its $(T,z)$ marginal is the original joint distribution given
$C=c$.  Conditional on the first endpoint $z$, the tube law is
\begin{equation}\label{eq:transverse-posterior-tube-law}
 \Prob(T\in\mathcal A\mid z,C=c)
   =\sum_{p:X_{s+b}(p)=z}
       \Prob(X=p\mid z,C=c)\Prob(T\in\mathcal A\mid X=p).
\end{equation}
The parent is determined by $z$, and $X$ determines all these
spatial labels.  Equation \eqref{eq:src1} therefore applies to
every inner law.  Use the slightly coarser rational angular scale
$L^{-b/3}=K^{-1/3}$, which covers the above cap for large $L$ and
has $b/3<u_0$.  For every fixed far pair the resulting angular
probability is at most
\[
 K^{-\kappa/3+o(1)}.
\]
For each first endpoint, \eqref{eq:transverse-partition-occupancy}
allows at most $K^{3\varepsilon+O(\eta)}$ same-open-cell second
endpoints.  Each good second incidence costs an additional
$K^{-1/2+O(\eta)}$ under its original $q_T$ law.  For each first $z$ this list is fixed before the tube is drawn.
In the integral over the original law of $T$ given $z,C$, apply the
bound on $q_T(z')$ only where $(T,z')$ is good, and apply the
angular cap bound to the resulting indicator.  Sum over the fixed
list of $z'$ and then over the first-endpoint law in
\eqref{eq:transverse-posterior-tube-law}.
The total far-pair contribution is at most
\begin{equation}\label{eq:transverse-far-pairs}
 K^{-1/2-\kappa/3+3\varepsilon+O(\eta)+o(1)}.
\end{equation}
No angular regularity after conditioning on the success event has
been used; success only restricts the pairs to which the original
mass bounds apply.

Divide \eqref{eq:transverse-near-pairs} and
\eqref{eq:transverse-far-pairs} by
\eqref{eq:transverse-pair-lower}.  The respective ratios are bounded
by
\[
 K^{-\varepsilon+O(\eta)},\qquad
 K^{-\kappa/3+2\varepsilon+O(\eta)+o(1)}.
\]
Choose $\varepsilon>0$ sufficiently small relative to $\kappa$,
then $\eta>0$ sufficiently small relative to $\varepsilon$.
Both ratios tend to zero, contradicting
\eqref{eq:transverse-pair-lower}.
\end{proof}

\begin{remark}[A grid proof of the wall count]
Retain the notation of the proof of Lemma~\ref{lem:transverse-no-jump}.
For comparison, the same wall count follows from a generic grid of
spacing $r$.  A component wholly contained inside a grid box costs
at most $O(D^3)$ boxes by the real algebraic component bound
\cite[Lemma~(C), p.~181]{W93}.
All other intersections are counted on grid faces.  On each
generic grid plane the curve meets $O(Dr^{-1}+D^2)$ squares:
intersections with grid lines cost $O(Dr^{-1})$, and components
wholly inside squares cost $O(D^2)$.  There are $O(r^{-1})$
relevant planes, giving
$O(Dr^{-2}+D^2r^{-1}+D^3)$.  A bounded enlargement covers near-wall
boxes.  The volume argument above already includes singular zero
sets and suffices for the estimate.
\end{remark}

\begin{proof}[Proof of Proposition~\ref{prop:transverse}]
A countersequence gives the weighted tube configuration by
Lemma~\ref{lem:transverse-tubes}.  Its limiting information profiles
satisfy \eqref{eq:src4}.  Lemma~\ref{lem:transverse-binary} makes
the common derivative binary, and
Lemma~\ref{lem:transverse-finite-jumps} supplies a positive-probability
full-to-empty block for any prescribed fixed small tolerance.
Lemma~\ref{lem:transverse-no-jump} rules out such a block.  Hence
no countersequence exists, which proves the existence of the fixed
power $c>0$ and the large-scale threshold.
\end{proof}
